
\begin{abstract}
The Vera C. Rubin Observatory will generate petabytes of data through the Legacy Survey of Space and Time (LSST) over the next decade, enabling discoveries across a broad range of astrophysical fields. Alongside these data products, Rubin maintains a large but heterogeneous collection of supporting documentation, including operational guides, technical notes, and scientific papers. Because this material is distributed across multiple platforms and formats, staff and scientists often struggle to efficiently locate accurate, up-to-date information. Many resources also reside on internal systems, limiting the ability of general-purpose language models to provide reliable answers to Rubin-specific questions. To address these challenges, we explore the use of Retrieval-Augmented Generation (RAG) to improve information discovery. We present a prototype RAG-based virtual assistant that delivers context-aware, factual, conversational access to Rubin’s vast and heterogenous documentation ecosystem. The system integrates material from multiple sources and enables semantic search through a conversational interface, using Weaviate for embeddings, LangChain for query orchestration, and an OpenAI GPT model as the LLM backend. By grounding responses in domain-specific knowledge, the assistant reduces hallucinations, improves accuracy, and demonstrates the potential of RAG to enhance access to distributed knowledge, streamline workflows, and support effective use of LSST data products.
\end{abstract}

